\documentclass[a4paper,12pt]{article}

\usepackage[a4paper,margin=1in]{geometry}
\usepackage{enumitem}
\usepackage{graphicx}
\usepackage{listings}
\usepackage{xcolor}
\usepackage{hyperref}

\setlength{\parindent}{0pt}
\setlength{\parskip}{0.6em}

% C Code Style
\lstdefinestyle{CStyle}{
    language=C,
    basicstyle=\ttfamily\footnotesize,
    keywordstyle=\color{blue},
    commentstyle=\color{green!40!black},
    stringstyle=\color{red!70!black},
    numbers=left,
    numberstyle=\tiny,
    stepnumber=1,
    breaklines=true,
    frame=single
}

\title{Assignment 1 Report: \texttt{spireslayer}\\
\large CMPE 230 Systems Programming -- Spring 2026}

\author{
Student 1 Name, Student 1 ID \\
Student 2 Name, Student 2 ID
}

\date{\today}

\begin{document}

\maketitle

\section*{Report Guidelines}

This report should explain \textbf{your implementation decisions}, not repeat the assignment handout.

Be specific. Instead of writing general statements about C or programming, describe the \textbf{actual design choices you made} in your program.

Useful writing tips:
\begin{itemize}[leftmargin=1.5em]
\item Refer to your data structures and helper functions.
\item Explain tricky cases or design tradeoffs.
\item Use short examples from your implementation or tests.
\end{itemize}

\section{Project Overview}

Briefly describe the overall structure of your program.

You may discuss:
\begin{itemize}[leftmargin=1.5em]
\item what the program does,
\item what the persistent state of the program is,
\item what design priorities shaped your implementation (simplicity, validation, modularity, debugging, etc.).
\end{itemize}

\section{Interpreter Design}

Explain the overall design of your interpreter.

\subsection{Input Handling and Parsing}

Describe how one input line is processed.

Possible points:
\begin{itemize}[leftmargin=1.5em]
\item how input lines are read,
\item how commands are recognized,
\item how the program checks whether a line fully matches the grammar,
\item how invalid commands are detected.
\end{itemize}

\subsection{State Representation}

Describe how the game state is stored in memory.

Examples:
\begin{itemize}[leftmargin=1.5em]
\item gold, HP, floor, current room,
\item cards and upgrades,
\item relics and potions,
\item codex effectiveness notes,
\item defeated counts and exhaust tags.
\end{itemize}

Explain why your chosen representation (arrays, structs, counters, etc.) works well for this assignment.

\subsection{Command Execution Flow}

Choose one or two representative commands and explain how they are executed.

Describe:
\begin{itemize}[leftmargin=1.5em]
\item which function processes the command,
\item what checks are performed,
\item what state changes occur,
\item what output is produced.
\end{itemize}

\section{Implementation Details}

Discuss parts of the implementation that required careful design or debugging.

Possible topics:
\begin{itemize}[leftmargin=1.5em]
\item whole-line validation,
\item sorting before output,
\item enforcing constraints (unique relics, potion limits, etc.),
\item combat or room transition logic,
\item handling malformed input.
\end{itemize}

Focus on a few important cases rather than describing the entire program.

\subsection{Representative Code Snippet}

Include a short excerpt (10--20 lines) from your implementation.

\begin{lstlisting}[style=CStyle]
/* Example snippet from your project */
\end{lstlisting}

Explain:
\begin{itemize}[leftmargin=1.5em]
\item what the code does,
\item why it is important,
\item what mistakes could occur here,
\item why your implementation works correctly.
\end{itemize}

\section{Testing and Validation}

Explain how you verified that your program works correctly.

\subsection{Testing Strategy}

Describe how you tested the program.

Examples:
\begin{itemize}[leftmargin=1.5em]
\item provided grader or checker,
\item custom input files,
\item manual interactive testing,
\item compiler warnings or debugging prints.
\end{itemize}

\subsection{Sample Interaction}

Include a short input/output example demonstrating several commands.

\subsection{Edge Cases}

List important edge cases you tested, for example:

\begin{itemize}[leftmargin=1.5em]
\item malformed commands,
\item extra tokens,
\item missing items,
\item duplicate relics,
\item full potion belt,
\item leading zeros in integers,
\item sorting behavior in queries.
\end{itemize}

\section{Discussion of Course Concepts}

Relate your implementation to concepts from the course.

\subsection{Data Representation and Memory Organization}

Explain how the program state is organized in memory and why this design works well.

\subsection{Strings, Arrays, and Pointers}

Discuss any non-trivial string processing, array traversal, or pointer usage in your solution.

\subsection{Memory Management Choices}

Explain your memory strategy.

\begin{itemize}[leftmargin=1.5em]
\item If you used dynamic memory, explain how allocations and frees are handled.
\item If you used static arrays, explain why they were sufficient.
\end{itemize}

\subsection{Robustness and Debugging}

Describe one debugging challenge you encountered and how you solved it.

\subsection{Limitations and Tradeoffs}

Describe any limitations in your design or implementation.

\section{Conclusion}

Summarize your implementation briefly.

Possible points:
\begin{itemize}[leftmargin=1.5em]
\item strongest part of the design,
\item improvements you would make with more time.
\end{itemize}

\section*{AI Usage Statement}

The report must be written by you.

If you used AI tools, document them here:
\begin{itemize}[leftmargin=1.5em]
\item tool name,
\item how it was used (e.g., grammar checking, code formatting, brainstorming, debugging help, etc.),
\item confirmation that the technical content is your own.
\end{itemize}

\end{document}